\chapter{Introduction}
\label{ch:introduction}
Searching a large embedding database requires decisions about both
information and work: how compactly to represent each vector, which vectors
to inspect, and how much precision to spend on them. A GPU makes many
operations parallel, but it does not make these decisions free. This thesis
studies their interaction in JHQ-GPU, a GPU implementation and experimental
analysis of JHQ's primary/residual quantisation hierarchy.

\section{Context and Practical Objectives}
\label{sec:1-1}
Approximate nearest-neighbour (ANN) search reduces query time by avoiding
exhaustive, full-precision distance evaluation. Graph methods such as CAGRA
restrict visits through navigation, while compressed IVF combines selective
list access with inexpensive scoring \cite{cagra,johnson_gpu,ivfrabitq}.
The tested embeddings have 768--3,072 coordinates and up to 17.78 million
vectors; Stella's raw FP32 array alone requires 67.81\,GiB
(\cref{sec:7-7}). Compact codes address capacity, but table construction,
lookups, selection and launch overhead determine their query-time benefit.

JHQ separates screening from refinement: compact primary codes score
candidates, and residuals refine a shortlist \cite{jhq}. This exposes two
decisions: how to execute the broad screen and how many survivors to refine.
The trained index can be reused while varying execution or budget, allowing
their effects to be tested separately. This thesis evaluates these decisions
through Recall@10, host-visible throughput/latency, resident memory and
construction cost, including configurations where a competitor or the
unmodified design is faster.


\section{Problem Definition and Scope}
\label{sec:1-2}
For a database $\mathcal X\subset\mathbb R^d$ and query $q$, the target is
the top-$k$ set under raw squared L2. The implemented system routes through
IVF, ranks the visited vectors by a primary reconstructed score
$\primarydist$, retains up to $c=\max\{k,\lceil\alpha k\rceil\}$ candidates,
and selects results using a residual-refined score $\refineddist$.
Both scores differ from raw-vector distance. The refined value additionally
contains a query-only constant, as defined in \cref{sec:2-3}.
Exactness of a selection procedure therefore concerns its computed score
and tie contract; final ANN quality remains an empirical question against GT.

The representation follows JHQ: a square orthogonal rotation, Cartesian
primary codes and residuals relative to primary reconstruction. This thesis
studies GPU training/encoding, physical placement, query execution and one
workload-level refinement multiplier. The multiplier is selected from
unlabelled query outputs against a larger-budget reference, then is intended
to be reused under a stable workload. It is not a per-query learned controller
or an absolute-recall certificate.

The main cross-system experiments use one RTX 5090, squared-L2 top-10 and
batches of 1,000 queries on four datasets. Six datasets support component
and budget studies, with separate large-index measurements. T5 adds
cross-system batches 8/32 alongside internal scheduling ablations
(\cref{sec:7-smallbatch}). The implemented compressed index
resides on one GPU; CPU model preparation and host-query transfers remain
part of its actual architecture. Multi-GPU distribution, online updates,
host-backed residual paging and billion-scale deployment are outside the
tested system.


\section{Three Linked Technical Challenges}
\label{sec:1-3}
\Cref{fig:1-challenges} connects three challenges to the mechanisms and
evidence used to address them.

First, the primary codebook has useful mathematical structure that a direct
GPU port can overlook. Each eight-bit subspace code indexes 256 distances.
Its Cartesian coordinates permit an exact 16+16 factorisation, but more
aggressive splitting adds lookups and arithmetic. The challenge is to exploit
separability without assuming that the smallest table gives the fastest
scan. Candidate-parallel word access and residual lane cooperation also need
different data mappings.

Second, a query's primary winners can be concentrated in any part of its
routed pool. Parallel groups cannot safely retain only a fraction of the
required global survivors. Exact local retention, a global merge and a
shared complete-score bound provide a principled route, but add buffers,
staging and synchronisation. A large batch can already supply enough blocks;
an imbalanced workload can still benefit from more groups. The challenge
couples correctness with a workload-dependent execution decision. Selecting
IVF lists remains a separate primitive from retaining database candidates.

Third, a fixed survivor multiplier can repeatedly refine irrelevant vectors,
while a small multiplier can remove a true neighbour too early. Agreement
with a large-budget reference is available without GT, but agreement does
not repair the reference's own errors or guarantee future-query recall.
Choosing a sample and budget also consumes time. The challenge is to evaluate
selection quality, calibration cost and reuse together.

\par\medskip\noindent\begin{minipage}{\linewidth}
\centering\captionsetup{hypcap=false}
\resizebox{\linewidth}{!}{\begin{tikzpicture}[
 font=\small\linespread{1}\selectfont,
 box/.style={draw=ThesisNavy!65,rounded corners=2pt,align=center,
 text width=3.5cm,minimum height=1.65cm,inner sep=4pt},
 arr/.style={-{Latex[length=2mm]},draw=ThesisNavy!80,thick}]
\node[font=\bfseries] at (0,1.25) {Challenge};
\node[font=\bfseries] at (4.6,1.25) {Mechanism};
\node[font=\bfseries] at (9.2,1.25) {Evidence};
\node[box] (a) at (0,0) {Primary codes\\Broad scan cost};
\node[box,fill=ThesisNavy!7] (b) at (4.6,0) {Cartesian tables\\Packed word scan};
\node[box] (c) at (9.2,0) {Table / load tests\\Score checks};
\node[box] (d) at (0,-2.2) {Uneven query work\\Global top-$c$};
\node[box,fill=ThesisNavy!7] (e) at (4.6,-2.2) {Grouped scan\\Shared bound\\Global merge};
\node[box] (f) at (9.2,-2.2) {Output checks\\Batch / skew tests};
\node[box] (g) at (0,-4.4) {Refinement effort\\Quality / reuse cost};
\node[box,fill=ThesisNavy!7] (h) at (4.6,-4.4) {Sampled agreement\\Workload budget};
\node[box] (i) at (9.2,-4.4) {Held-out quality\\Calibration cost};
\draw[arr] (a)--(b);\draw[arr] (b)--(c);
\draw[arr] (d)--(e);\draw[arr] (e)--(f);
\draw[arr] (g)--(h);\draw[arr] (h)--(i);
\end{tikzpicture}
}
\captionof{figure}[Problem, mechanism and evidence]{Problem, mechanism and evidence. The arrows organise
research questions; they do not imply that independently measured speedups
multiply or that an unlabelled agreement rule guarantees GT recall.}
\label{fig:1-challenges}
\end{minipage}\par\medskip


\section{Research Questions and Falsifiable Claims}
\label{sec:1-4}
\paragraph{RQ1: Construction and capacity.}
How can JHQ's hierarchy be trained, encoded and arranged as a searchable
GPU-resident index, and what construction cost and tested capacity result?
The evidence is the physical layout, train/add records and resident-memory
account (\cref{sec:7-7}). An unsuccessful build or an allocation beyond the
device budget would limit the capacity claim; a resident sample alone does
not establish construction peak.

\paragraph{RQ2: Correct and efficient execution.}
Which representation-aware scoring, retention and scheduling choices
improve GPU search while respecting the specified score and tie semantics?
Paired LUT, packed-load, refinement and scan experiments connect to reference
checks and batch sensitivity (\cref{sec:7-3,sec:7-4,sec:7-6}). A non-tie
selection error would challenge the correctness contract; consistent
slowdowns would challenge the proposed execution choice.

\paragraph{RQ3: Budget quality and total cost.}
Can unlabelled workload agreement choose a useful survivor budget, and under
what sampling and reuse conditions is calibration worthwhile? Disjoint query
evaluation compares fixed budgets, sampled/FULL rules, enumeration and a
GT-labelled oracle (\cref{sec:7-5}). Loss crossings challenge reliability;
no positive saving after calibration challenges the net-cost benefit.
The oracle is an evaluation control, not deployable label-free logic.

\paragraph{RQ4: Competitive operating regions.}
At which measured recall targets does the complete system compete with
modern GPU ANN alternatives under a common host-to-host boundary?
The answer uses independently tuned E1 frontiers and explicit memory
definitions (\cref{sec:7-2}). Int8 graph competitors, high-recall reversals
and absent coverage are part of the answer. The question does not seek a
universal ranking across hardware, memory budgets or request sizes.

Chapter 3 turns these questions into controlled comparisons. Together they
separate what was built, whether its execution is correct, what effort a
workload requires, and whether the resulting system has useful operating
regions.


\section{Contributions and Research Ownership}
\label{sec:1-5}
\Cref{tab:1-contributions} states three linked project contributions and
their inherited foundations. The core contribution is the adaptation and
evaluation of a particular representation--execution--budget combination.
JHQ, Cartesian quantisation, IVF, bitonic/radix selection, merging and CUDA
Graph execution retain their prior-art status.

\par\medskip\noindent\begin{minipage}{\linewidth}
\centering\captionsetup{hypcap=false}
\captionof{table}[Technical contributions and evidence boundaries]{Technical contributions and their evidence boundaries.}
\label{tab:1-contributions}
\footnotesize\setlength{\tabcolsep}{4pt}
\begin{tabularx}{\linewidth}{@{}l*{3}{>{\raggedright\arraybackslash}X}@{}}
\toprule
 & Inherited foundation & Project contribution & Evidence \\
\midrule
C1 & JHQ rotation, codebooks and hierarchy & GPU training/encoding,
bounded batch construction and list-ordered packed representation
& Search of 17.78M vectors at 20.44\,GiB resident; repeated train+add
10.54--10.65\,s; sampled construction peak 22.58\,GiB \\
C2 & Cartesian additivity and established GPU primitives & Integrated
factorised scoring, complete-key grouped retention/bounds, lane refinement
and execution scheduling & Two-group LUTs lead in 22/24 cells;
scheduled execution $2.727\times$ at batch 8; explicit correctness scope \\
C3 & Candidate expansion and output-agreement ideas & Workload-level
budget rule, conditional selection analysis, held-out controls and cost/risk
evaluation & FULL meets the relative-loss target in 12/12 cells;
smaller samples expose risk; four cells directly measure calibrated reuse \\
\bottomrule
\end{tabularx}
\end{minipage}\par\medskip

On the full-batch E1 grid near recall 0.97, JHQ leads in three datasets and is
near parity with RaBitQ on arXiv under the interpolated comparison. Other
recall regions favour competitors. FULL's maximum observed held-out loss is
0.0006 relative to the best tested arm, while $S=128$ reaches 40/64 loss
crossings in one cell. These are empirical quality results; ten FULL cells
retain conditional nestedness coverage, while all twelve FULL
finite-grid choices now match enumeration (\cref{sec:7-budget-trace}).
The OpenAI E1 results use TF32 GEMMs; a separate FP32 sensitivity
retains smaller favourable margins (\cref{sec:7-precision}).

This COMP8800 project is presented by Rui Luo under Dr Mengxuan Zhang's
supervision. The coauthored ADC manuscript names Rui Luo and Mengxuan Zhang
and shares the JHQ-GPU implementation and measurement artifacts
\cite{adc_artifact}. The thesis
expands the correctness contracts, newer experiment analysis and treatment
of failed designs; historical manuscript numbers retain their own versions.
Technical contribution labels describe the project scope. The acknowledgements
identify joint work and AI assistance.
The thesis author is responsible for the submitted claims and attribution.


\section{Thesis Roadmap}
\label{sec:1-6}
Chapters 2 and 3 establish the representation, alternatives and measurement
design. Chapter 4 follows training and physical index construction;
Chapter 5 develops GPU scoring and selection contracts. Chapter 6 analyses
workload-level refinement calibration. Chapter 7 evaluates construction,
execution, budget quality and competitive operating regions. Chapter 8
interprets those findings through negative results, validity limits and
testable extensions. Chapter 9 answers the research questions and states
the resulting systems principle. Appendix A maps claims to implementation
and measurement artifacts.
